\subsection{Checklist: Open access to data and code}
\label{app:checklist_open_data_code}

\paragraph{Question.}
Does the paper provide open access to the data and code needed to reproduce the
main experimental results?

\paragraph{Answer (Yes/No/NA).}
Yes.

\paragraph{Justification.}
The empirical setting uses generated dynamical-system observations rather than a
restricted or proprietary dataset. The controlled multibasin benchmark consists
of procedurally specified two-dimensional systems with known attractor geometry,
and the Dysts benchmark uses public chaotic-flow systems from the external
\texttt{dysts} package. The released artifact will include the implementation
used to generate observations, train sparse and dense Koopman autoencoders,
evaluate held-out rollouts, compute support diagnostics, and build the
paper-facing tables and figures. Basin labels, basin counts, and attractor
coordinates are included only as benchmark metadata for evaluation and
diagnostics; the training code does not require them for sparse-support
selection.

\paragraph{Exactly what will be released and included.}
\begin{itemize}
  \item The core source package, including model definitions, sparse encoders,
  Koopman transition variants, configuration presets, dynamical-system
  environments, benchmark adapters, and evaluation utilities.
  \item Command-line tools and batch scripts for training, checkpoint
  evaluation, benchmark task construction, result collection, statistical
  summaries, and figure/table generation.
  \item Environment and reproducibility metadata, including
  \texttt{pyproject.toml}, \texttt{uv.lock}, the project README, and the
  package entry points for training and evaluation.
  \item The paper source and paper-facing generated artifacts needed to inspect
  the reported claims, including the LaTeX tables, figure files, aggregate
  summaries, and collected CSV/JSON metric files used to produce the main
  forecasting and support-diagnostic displays.
  \item Benchmark definitions for the 15 controlled multibasin systems and the
  10 Dysts \(dt{\times}30\) forecasting systems, including system keys,
  dimensions, stored-step conventions, and the evaluation-only metadata needed
  to score basin-support alignment on the controlled benchmark.
  \item Instructions for regenerating trajectories, rerunning the training and
  evaluation pipeline, and rebuilding the paper summaries from released
  artifacts.
\end{itemize}

\paragraph{Data generation and access.}
For the controlled multibasin benchmark, there is no separate private dataset.
Training, validation, and test windows are generated from released vector-field
implementations by applying the stored observation map with fourth-order
Runge--Kutta integration. The learner observes stored states only; it does not
observe continuous-time vector fields, integration substeps, basin labels, or
trajectory-to-basin assignments during training. Evaluation-only basin metadata
is used after training to define diagnostic slices, controlled transfer pairs,
and support--basin agreement scores.

For the Dysts benchmark, trajectories are generated from public Dysts systems
through the released adapter code and the dependency version resolved by the
project environment. The paper-facing setup uses standardized coordinates,
native Dysts trajectory generation, deterministic train/validation/test cache
split namespaces, 200 cached trajectories, 30,000 stored observations per
trajectory, and a 2,000-step warm-up before cached observations are used. The
released code can rebuild these caches from the public Dysts right-hand sides;
the cache files themselves are generated artifacts rather than required external
data.

\paragraph{Anonymization.}
The experiments do not use human-subject data, personal data, private text,
images, medical records, user logs, or other sensitive data. The generated
states, labels, checkpoints, and metrics contain only synthetic dynamical-system
values and model outputs. During anonymous review, the release will remove or
neutralize author-identifying metadata, local usernames, private scratch paths,
and cluster-specific job records that are not needed for reproduction. The final
public release can restore standard project attribution and licensing.

\paragraph{What is not released.}
The release will not include private cluster scratch directories, raw SLURM
stdout/stderr logs, queue bookkeeping, editor/agent state, exploratory planning
notes, failed pilot runs, or the full \texttt{runs/} tree of transient training
outputs. Large generated Dysts cache tensors and exhaustive checkpoint files for
exploratory sweeps are not required for access to the benchmark data, because
they can be regenerated from the released code and public dependencies. The
paper-facing aggregate CSV/JSON summaries, tables, figures, and scripts needed
to audit the reported results will be included.
